\documentclass{beamer}
\usepackage[utf8]{inputenc}

\usetheme{Madrid}
\usecolortheme{default}
\usepackage{amsmath,amssymb,amsfonts,amsthm}
\usepackage{txfonts}
\usepackage{tkz-euclide}
\usepackage{listings}
\usepackage{adjustbox}
\usepackage{array}
\usepackage{tabularx}
\usepackage{gvv}
\usepackage{lmodern}
\usepackage{circuitikz}
\usepackage{tikz}
\usepackage{graphicx}

\setbeamertemplate{page number in head/foot}[totalframumber]

\usepackage{tcolorbox}
\tcbuselibrary{minted,breakable,xparse,skins}

\definecolor{bg}{gray}{0.95}
\lstset{
    language=C,
    basicstyle=\ttfamily\small,
    keywordstyle=\color{blue},
    stringstyle=\color{orange},
    commentstyle=\color{green!60!black},
    numbers=left,
    numberstyle=\tiny\color{gray},
    breaklines=true,
    showstringspaces=false,
}

%------------------------------------------------------------
\title{5.4.6}
\date{September 30,2025}
\author{EE25BTECH11011 - Navya Banavathu}

\begin{document}

\frame{\titlepage}

\begin{frame}{Question}
 Using elementary transformations, find the inverse of the following matrix. 
\begin{align*}
    \myvec{10&&-2\\-5&&1}
\end{align*}

\end{frame}

\begin{frame}{Solution}
We employ the Gauss--Jordan method,
\begin{align}
\left(\begin{array}{cc|cc}
10 & -2 & 1 & 0 \\
-5 & 1 & 0 & 1
\end{array}\right)
&\xrightarrow{R_1 \leftarrow \tfrac{1}{10}R_1}
\left(\begin{array}{cc|cc}
1 & -\tfrac{1}{5} & \tfrac{1}{10} & 0 \\
-5 & 1 & 0 & 1
\end{array}\right) \\[1em]
&\xrightarrow{R_2 \leftarrow R_2+5R_1}
\left(\begin{array}{cc|cc}
1 & -\tfrac{1}{5} & \tfrac{1}{10} & 0 \\
0 & 0 & \tfrac{1}{2} & 1
\end{array}\right)
\end{align}

At this stage, the left-hand side is not the identity matrix. \\
Hence the given matrix is not invertible.

\end{frame}


\begin{frame}[fragile]{ Code (C)}
\begin{block}{Part 1}
\begin{lstlisting}[language=C, basicstyle=\ttfamily\small, keywordstyle=\color{blue}, frame=single]
#include <stdio.h>

void compute_inverse(double* inv) {
    double a = 10, b = -2;
    double c = -5, d = 1;
    double det = a*d - b*c;

    if(det == 0) {
        inv[0] = inv[1] = inv[2] = inv[3] = 0;
    } else {
        inv[0] = d/det;   // [0,0]
        inv[1] = -b/det;  // [0,1]
        inv[2] = -c/det;  // [1,0]
        inv[3] = a/det;   // [1,1]
    }
}
\end{lstlisting}
\end{block}
\end{frame}


\begin{frame}[fragile]{Python Code}
\begin{block}{Part 1}
\begin{lstlisting}[language=Python, basicstyle=\ttfamily\small, keywordstyle=\color{blue}, frame=single]
import ctypes
import numpy as np
import os

# Load the shared library
if os.name == "nt":
    lib = ctypes.CDLL("./code9.dll")
else:
    lib = ctypes.CDLL("./code9.so")

# Create a 2x2 NumPy array to hold the inverse
inv_matrix = np.zeros((2,2), dtype=np.double)

# Set the argument type for ctypes
lib.compute_inverse.argtypes = [ctypes.POINTER(ctypes.c_double)]
lib.compute_inverse.restype = None
\end{lstlisting}
\end{block}
\end{frame}


\begin{frame}[fragile]{Python Code}
\begin{block}{Part 2}
\begin{lstlisting}[language=Python, basicstyle=\ttfamily\small, keywordstyle=\color{blue}, frame=single]
# Call the C function with pointer to NumPy array
lib.compute_inverse(inv_matrix.ctypes.data_as(ctypes.POINTER(ctypes.c_double)))

# Check if matrix is invertible
if np.all(inv_matrix == 0):
    print("Matrix is not invertible.")
else:
    print("Matrix is invertible.")
    print("Inverse matrix:\n", inv_matrix)


\end{lstlisting}
\end{block}
\end{frame}
\end{document}

